\section{The Jacobian as an abelian variety}
\label{mk-chap:Jacobian}

Let \(C\) be a smooth proper geometrically irreducible curve over a field
\(k\), with genus \(g(C)\).  The Jacobian is constructed uniformly as the
identity component \(\Pic^0_{C/k}\) of the Picard scheme.  It is an abelian
variety of dimension \(g(C)\); when \(g(C)=0\), it is \(\Spec k\).

\section{Albanese objects}

\begin{definition}[Albanese universal property of a pointed curve]\label{mk-def:IsAlbanese}
  \lean{AlgebraicGeometry.IsAlbanese}
  \leanok
  \dcref{custom}

  Let \(P \in C(k)\).  An abelian variety \(J\) together with a pointed
  morphism \(\iota_P \colon C \to J\) is an \emph{Albanese object} for
  \((C,P)\) if every pointed morphism \(f \colon C \to A\) to an abelian
  variety factors uniquely through \(\iota_P\): there is a unique morphism of
  \(k\)-schemes \(g \colon J \to A\) such that
  \[
    P \circ \iota_P = \eta_J,
    \qquad
    f = \iota_P \circ g.
  \]
\end{definition}

\begin{remark}
  \label{mk-rem:IsAlbanese_typeclasses}
  \dcref{custom}
  The phrase ``abelian variety'' includes the group structure, properness,
  smoothness, and geometric irreducibility of both \(J\) and every target
  \(A\).  These are hypotheses of the universal property, rather than extra
  conclusions inside it.
\end{remark}

\subsection{The universal morphism}

\begin{definition}[The Albanese universal morphism]\label{mk-def:IsAlbanese_ofCurve}
  \lean{AlgebraicGeometry.IsAlbanese.ofCurve}
  \leanok
  \dcref{custom}
  \uses{mk-def:IsAlbanese}

  An Albanese structure on \((C,P,J)\) determines a distinguished morphism
  \(\iota_P \colon C \to J\), called its universal morphism.
\end{definition}

\begin{lemma}[Pointed condition]\label{mk-lem:IsAlbanese_comp_ofCurve}
  \lean{AlgebraicGeometry.IsAlbanese.comp_ofCurve}
  \leanok
  \dcref{custom}
  \uses{mk-def:IsAlbanese, mk-def:IsAlbanese_ofCurve}

  The universal morphism satisfies \(P \circ \iota_P = \eta_J\).
\end{lemma}

\begin{lemma}[Universal factorisation]\label{mk-lem:IsAlbanese_exists_unique_ofCurve_comp}
  \lean{AlgebraicGeometry.IsAlbanese.exists_unique_ofCurve_comp}
  \leanok
  \dcref{custom}
  \uses{mk-def:IsAlbanese, mk-def:IsAlbanese_ofCurve}

  Every pointed morphism \(f \colon C \to A\) to an abelian variety factors
  uniquely as \(f=\iota_P\circ g\) through a morphism of \(k\)-schemes
  \(g \colon J \to A\).
\end{lemma}

\begin{theorem}[Uniqueness of the Albanese object]\label{mk-thm:IsAlbanese_unique}
  \lean{AlgebraicGeometry.IsAlbanese.unique}
  \leanok
  \dcref{custom}
  \uses{mk-def:IsAlbanese, mk-def:IsAlbanese_ofCurve}

  If \((J_1,\iota_1)\) and \((J_2,\iota_2)\) are Albanese objects for
  \((C,P)\), there is a unique morphism \(e \colon J_1\to J_2\) satisfying
  \(\iota_2=\iota_1\circ e\).
\end{theorem}

\begin{proof}
  \leanokThe universal property of \(J_1\) applied to \(\iota_2\) gives
  \(e \colon J_1\to J_2\); that of \(J_2\) applied to \(\iota_1\) gives
  \(e' \colon J_2\to J_1\).  Both \(e\circ e'\) and the identity of \(J_1\)
  factor \(\iota_1\) through itself, so uniqueness gives
  \(e\circ e'=\mathrm{id}_{J_1}\); similarly
  \(e'\circ e=\mathrm{id}_{J_2}\).  The original universal property gives the
  required uniqueness of \(e\).
\end{proof}

\begin{remark}
  \label{mk-rem:IsAlbanese_unique_iso}
  \dcref{custom}
  Consequently the Albanese object is unique up to a unique isomorphism
  compatible with the universal morphism.
\end{remark}

\section{A uniform Jacobian witness}

\begin{definition}[Jacobian witness]\label{mk-def:JacobianWitness}
  \lean{AlgebraicGeometry.JacobianWitness}
  \leanok
  \dcref{custom}
  \uses{mk-def:IsAlbanese, mk-def:genus}

  A \emph{Jacobian witness} for \(C\) consists of a single \(k\)-scheme \(J\),
  an abelian-variety structure on \(J\), a proof that \(J\) is smooth of
  relative dimension \(g(C)\), and, for every \(P\in C(k)\), an Albanese
  structure on \((C,P,J)\).
\end{definition}

\begin{remark}
  \label{mk-rem:JacobianWitness_quantifier_order}
  \dcref{custom}
  The same \(J\) works for every marked point.  Replacing \(P\) by another
  rational point translates the universal morphism by an automorphism of
  \(J\), so it does not change the underlying abelian variety.
\end{remark}

\begin{remark}
  \label{mk-rem:JacobianWitness_smooth_redundancy}
  \dcref{custom}
  Smoothness is implied by smoothness of relative dimension \(g(C)\).  It is
  nevertheless useful to record both the abelian-variety property and the
  dimension assertion as parts of the witness.
\end{remark}

\subsection{The Picard witness \(J=\Pic^0_{C/k}\)}
\label{mk-sec:picardJacobianWitness}

\begin{theorem}[A Jacobian witness for a pointed curve]
  \label{mk-def:picardJacobianWitnessOfHasRationalPoint}
  \lean{AlgebraicGeometry.picardJacobianWitnessOfHasRationalPoint}
  \dcref{custom}
  \uses{mk-def:JacobianWitness, mk-def:genus, mk-thm:fga_pic_representability,
    mk-thm:pic0_isAbelianVariety, mk-thm:pic_zero_dimension_equals_genus,
    mk-thm:albanese_universal_property_northstar, mk-def:has_rational_point}

  Let \(C\) be a smooth proper geometrically irreducible curve over \(k\)
  possessing a \(k\)-rational point. Then the identity component
  \(\Pic^0_{C/k}\) carries a Jacobian witness for \(C\).
\end{theorem}

\begin{proof}
  \uses{mk-lem:curve_geometrically_integral, mk-def:inst_has_pic_scheme,
    mk-lem:pic0_relative_dimension_genus, mk-lem:pic0_isAlbanese_all_points}
  Representability of the relative Picard functor
  (\ref{mk-def:inst_has_pic_scheme}, whose hypothesis is the rational point) gives
  \(\Pic_{C/k}\), and its identity component \(\Pic^0_{C/k}\) is a group scheme,
  smooth, proper and geometrically irreducible by
  \ref{mk-thm:pic0_isAbelianVariety}. Those four data are the first four fields of
  the witness; of the four, \ref{mk-thm:pic0_smooth} and \ref{mk-thm:pic0_proper} are
  the ones whose own proofs are still open.

  Three further steps are needed, and each is isolated as a lemma because it does
  not follow from the four above.

  First, the hypotheses must be reconciled: \(\Pic^0_{C/k}\) is constructed for a
  geometrically \emph{integral} curve, whereas a witness is required for a
  geometrically irreducible one. Geometric integrality is automatic
  (\ref{mk-lem:curve_geometrically_integral}).

  Second, the smoothness supplied by \ref{mk-thm:pic0_isAbelianVariety} carries no
  relative dimension, while the witness records smoothness at the specific
  relative dimension \(g(C)\); the tangent-space identification
  \(T_0\Pic^0_{C/k}\cong H^1(C,\mathcal O_C)\) of
  \ref{mk-thm:pic_zero_dimension_equals_genus} supplies the dimension, and
  \ref{mk-lem:pic0_relative_dimension_genus} is the resulting refinement.

  Third, the Albanese structure is required for \emph{every} \(P\in C(k)\) and for
  every genus, whereas \ref{mk-thm:albanese_universal_property_northstar} is proved
  over an algebraically closed base field and for \(g(C) > 0\); the extension is
  \ref{mk-lem:pic0_isAlbanese_all_points}. Together these seven data form the
  required witness.
\end{proof}

\begin{remark}[Pointed and unpointed witnesses]
  \label{mk-rem:pointed_witness_scope}
  \dcref{custom}
  \uses{mk-def:picardJacobianWitnessOfHasRationalPoint, mk-lem:curve_hypothesis_gap,
    mk-def:inst_has_pic_scheme, mk-thm:fga_pic_representability_etale}
  The pointed construction \ref{mk-def:picardJacobianWitnessOfHasRationalPoint}
  applies when \(C(k)\) is nonempty, using the \(\Pic^\sharp\)-representability
  of \ref{mk-def:inst_has_pic_scheme}.  Over an algebraically closed field this
  hypothesis follows from \ref{mk-lem:curve_rational_point_algClosed}.
  The unpointed witness \ref{mk-def:picardJacobianWitness} instead uses the
  etale-sheafified Picard functor of \ref{mk-thm:fga_pic_representability_etale}.
\end{remark}

\begin{theorem}[The Picard identity component is a Jacobian witness]
  \label{mk-def:picardJacobianWitness}
  \lean{AlgebraicGeometry.picardJacobianWitness}
  \dcref{custom}
  \uses{mk-def:JacobianWitness, mk-def:genus, mk-thm:fga_pic_representability_etale,
    mk-thm:pic0_isAbelianVariety, mk-thm:pic_zero_dimension_equals_genus,
    mk-thm:albanese_universal_property_northstar}

  Let \(k\) be an arbitrary field and \(C/k\) a smooth proper geometrically
  irreducible curve --- with \emph{no} hypothesis on \(C(k)\). Then the identity
  component \(\Pic^0_{C/k}\) carries a Jacobian witness for \(C\).
\end{theorem}

\begin{proof}
  \uses{mk-thm:fga_pic_representability_etale, mk-lem:pic0_relative_dimension_genus,
    mk-lem:pic0_isAlbanese_all_points, mk-rem:pointed_witness_scope,
    mk-rem:rational_point_scope}
  Take \(J = \Pic^0_{C/k}\), the identity component of the scheme representing the
  \'etale-sheafified relative Picard functor
  (\ref{mk-thm:fga_pic_representability_etale}). Its group-object structure, geometric
  irreducibility and local finiteness follow from that representability with no
  further hypothesis, by Yoneda transport and Kleiman \S 5 Lem.~\texttt{lem:agps};
  smoothness and properness are \ref{mk-thm:pic0_smooth} and \ref{mk-thm:pic0_proper};
  the relative dimension is \ref{mk-lem:pic0_relative_dimension_genus}; and the
  Albanese property for every marked point is
  \ref{mk-lem:pic0_isAlbanese_all_points}.

  No rational point is used, and none is available in general: see
  \ref{mk-rem:rational_point_scope} for why this rather than the pointed statement is
  the theorem being claimed, and \ref{mk-rem:pointed_witness_scope} for the conditional
  milestones kept beside it.
\end{proof}

\begin{lemma}[A smooth geometrically irreducible curve is geometrically integral]
  \label{mk-lem:curve_geometrically_integral}
  \lean{AlgebraicGeometry.geometricallyIntegral_of_curve}
  \leanok
  \dcref{custom}
  Let \(C\) be a smooth proper geometrically irreducible curve over a field \(k\).
  Then \(C\) is geometrically integral over \(k\).
\end{lemma}

\begin{proof}
  \leanok
  \uses{mk-lem:smooth_geometrically_reduced}
  A smooth morphism is geometrically reduced
  (\ref{mk-lem:smooth_geometrically_reduced}), and a morphism that is both
  geometrically reduced and geometrically irreducible is geometrically integral:
  for every field extension \(K/k\) the base change \(C_K\) is then both reduced
  and irreducible, hence integral.
\end{proof}

\begin{lemma}[A smooth morphism is geometrically reduced]
  \label{mk-lem:smooth_geometrically_reduced}
  \lean{AlgebraicGeometry.SmoothOfRelativeDimension.geometricallyReduced}
  \leanok
  \dcref{custom}
  A morphism of schemes \(f\colon X\to Y\) that is smooth of some relative
  dimension is geometrically reduced: for every field \(K\) and every morphism
  \(\Spec K\to Y\), the base change \(X\times_Y\Spec K\) is a reduced scheme.
\end{lemma}

\begin{proof}
  \leanok
  \uses{mk-lem:standard_smooth_over_domain_reduced}
  Smoothness is stable under base change, so it suffices to show that a scheme
  \(X\) smooth over a field \(K\) is reduced. Reducedness is local, and \(X\) is
  covered by affine opens \(\Spec S\) with \(S\) standard smooth over \(K\); each
  is reduced by \ref{mk-lem:standard_smooth_over_domain_reduced} applied to the
  domain \(K\).
\end{proof}

\begin{lemma}[A standard smooth algebra over a domain is reduced]
  \label{mk-lem:standard_smooth_over_domain_reduced}
  \lean{Algebra.IsStandardSmooth.isReduced_of_isDomain}
  \leanok
  \dcref{custom}
  Let \(B\) be a domain and \(S\) a standard smooth \(B\)-algebra. Then \(S\) is
  reduced.
\end{lemma}

\begin{proof}
  \leanok
  By the structure theorem for standard smooth algebras, \(S\) is étale over a
  polynomial ring \(P=B[x_1,\dots,x_n]\), which is again a domain. So it suffices
  to prove that a flat, formally unramified \(P\)-algebra \(S\) essentially of
  finite type over a domain \(P\) is reduced, étale algebras being flat, formally
  unramified and essentially of finite type.

  Let \(F=\operatorname{Frac}(P)\). The algebra \(F\otimes_P S\) is formally
  unramified and essentially of finite type over the field \(F\), hence reduced.
  Since \(S\) is flat over \(P\) and \(P\to F\) is injective, the induced map
  \(S\cong P\otimes_P S\to F\otimes_P S\) is injective, so \(S\) embeds into a
  reduced ring and is therefore reduced.
\end{proof}

\begin{remark}[The rational point is not a hypothesis of the Jacobian]
  \label{mk-lem:curve_hypothesis_gap}
  \dcref{custom}
  \uses{mk-def:has_rational_point, mk-thm:fga_pic_representability_etale}
  A smooth proper geometrically integral curve \(C/k\) need \emph{not} have a
  \(k\)-rational point: a conic over \(\mathbb{Q}\) with no rational point is a
  genus-\(0\) example, and there are genus-\(2\) curves over \(\mathbb{Q}\) with
  no rational point. Over an algebraically closed field a rational point always
  exists (\ref{mk-lem:curve_rational_point_algClosed}).

  The etale formulation uses no section, while the pointed formulation retains
  \(C(k)\neq\emptyset\) as an explicit hypothesis.
\end{remark}

\begin{lemma}[A curve over an algebraically closed field has a rational point]
  \label{mk-lem:curve_rational_point_algClosed}
  \lean{AlgebraicGeometry.hasRationalPoint_of_curve_of_isAlgClosed}
  \leanok
  \dcref{custom}
  \uses{mk-def:has_rational_point}
  Let \(C\) be a smooth proper geometrically irreducible curve over an algebraically
  closed field \(\bar k\). Then \(C\) has a \(\bar k\)-rational point.
\end{lemma}

\begin{proof}
  \leanok
  Geometric irreducibility over the one-point base \(\Spec \bar k\) makes the
  underlying space of \(C\) irreducible, in particular nonempty, and smoothness makes
  \(C \to \Spec \bar k\) locally of finite type, so that space is a Jacobson space.
  A nonempty Jacobson space has a closed point \(x\). Over an algebraically closed
  field the residue field at a closed point of a scheme locally of finite type is
  \(\bar k\) itself, so \(x\) is the image of a section of \(C \to \Spec \bar k\),
  which is the required rational point.
\end{proof}

\begin{remark}[The scope of the headline: an arbitrary base field]
  \label{mk-rem:rational_point_scope}
  \dcref{custom}
  \uses{mk-lem:curve_hypothesis_gap, mk-lem:curve_rational_point_algClosed,
    mk-thm:fga_pic_representability_etale}
  The Jacobian is defined over an arbitrary field from the etale sheafification
  \(\Pic_{(C/k)\et}\) and the five inputs
  \ref{mk-thm:fga_pic_representability_etale}, \ref{mk-thm:pic0_smooth},
  \ref{mk-thm:pic0_proper}, \ref{mk-lem:pic0_relative_dimension_genus}, and
  \ref{mk-lem:pic0_isAlbanese_all_points}.  The pointed and algebraically closed
  witnesses remain useful special cases.
\end{remark}

\begin{theorem}[A Jacobian witness over an algebraically closed field]
  \label{mk-def:picardJacobianWitnessOfIsAlgClosed}
  \lean{AlgebraicGeometry.picardJacobianWitnessOfIsAlgClosed}
  \dcref{custom}

  \uses{mk-def:JacobianWitness, mk-def:picardJacobianWitness,
    mk-lem:curve_rational_point_algClosed}
  Over an algebraically closed field \(\bar k\), the identity component
  \(\Pic^0_{C/\bar k}\) carries a Jacobian witness for every smooth proper
  geometrically irreducible curve \(C\), with no rational-point hypothesis.
\end{theorem}

\begin{proof}
  \uses{mk-def:picardJacobianWitnessOfHasRationalPoint,
    mk-lem:curve_rational_point_algClosed,
    mk-lem:pic0_relative_dimension_genus, mk-lem:pic0_isAlbanese_all_points,
    mk-rem:rational_point_scope, mk-def:inst_has_pic_scheme}
  Apply \ref{mk-def:picardJacobianWitnessOfHasRationalPoint}, with the rational point
  supplied by \ref{mk-lem:curve_rational_point_algClosed} instead of assumed. Every
  other field of the witness is unchanged, so the open obligations are
  representability (\ref{mk-def:inst_has_pic_scheme}), \ref{mk-thm:pic0_smooth},
  \ref{mk-thm:pic0_proper}, \ref{mk-lem:pic0_relative_dimension_genus} and
  \ref{mk-lem:pic0_isAlbanese_all_points}, all of them true statements
  (\ref{mk-rem:rational_point_scope}).
\end{proof}

\begin{lemma}[Relative dimension of the identity component]
  \label{mk-lem:pic0_relative_dimension_genus}
  \lean{AlgebraicGeometry.smoothOfRelativeDimension_genus_pic0}
  \dcref{custom}

  \uses{mk-def:genus, mk-thm:pic_zero_dimension_equals_genus, mk-thm:pic0_isAbelianVariety}
  For a smooth proper geometrically integral curve \(C/k\) whose Picard functor is
  representable, \(\Pic^0_{C/k}\) is smooth over \(k\) of relative dimension
  \(g(C)\).
\end{lemma}

\begin{proof}
  \(\Pic^0_{C/k}\) is a smooth \(k\)-group scheme by
  \ref{mk-thm:pic0_isAbelianVariety}. For a smooth group scheme over a field the
  relative dimension is constant and equals the dimension of the tangent space at
  the identity, which is \(H^1(C,\mathcal O_C)\) by
  \ref{mk-thm:pic_zero_dimension_equals_genus}. Its \(k\)-dimension is \(g(C)\) by
  the definition of the genus (\ref{mk-def:genus}).
\end{proof}

\begin{lemma}[Tangent space of the identity component has dimension the genus]
  \label{mk-lem:pic0_tangent_dimension_genus}
  \lean{AlgebraicGeometry.finrank_tangentSpace_pic0_eq_genus}
  \leanok
  \dcref{custom}
  \uses{mk-def:genus, mk-thm:pic_zero_dimension_equals_genus}
  For a smooth proper geometrically integral curve \(C/k\) whose Picard functor is
  representable, the Zariski tangent space of \(\Pic^0_{C/k}\) at the identity has
  dimension \(g(C)\) over the residue field of the identity point.
\end{lemma}

\begin{proof}
  \uses{mk-thm:pic_zero_dimension_equals_genus}
  The tangent space at the identity is \(H^1(C,\mathcal O_C)\)
  (\ref{mk-thm:pic_zero_dimension_equals_genus}), whose \(k\)-dimension is \(g(C)\) by
  the definition of the genus (\ref{mk-def:genus}).
\end{proof}

\begin{remark}[What separates the tangent dimension from the relative dimension]
  \label{mk-rem:tangent_versus_relative_dimension}
  \dcref{custom}
  \uses{mk-lem:pic0_tangent_dimension_genus, mk-lem:pic0_relative_dimension_genus}
  \ref{mk-lem:pic0_tangent_dimension_genus} is not
  \ref{mk-lem:pic0_relative_dimension_genus}. Two further ingredients separate them:
  smoothness of \(\Pic^0_{C/k}\) over \(k\), and the identification of the
  tangent-space dimension with the relative dimension of a smooth morphism. For the
  second, the relative dimension of a smooth morphism is computed by the rank of the
  module of Kähler differentials \(\Omega\), so the step is the duality, over a field,
  between that rank and the dimension of the tangent space at a rational point ---
  together with the passage from that affine-local statement to a cover of
  \(\Pic^0_{C/k}\). Translation by a rational point makes the dimension independent of
  the chosen point, which is what lets a computation at the identity determine the
  relative dimension everywhere.
\end{remark}

\begin{lemma}[Albanese property at every point, in every genus]
  \label{mk-lem:pic0_isAlbanese_all_points}
  \lean{AlgebraicGeometry.isAlbanese_pic0}
  \dcref{custom}

  \uses{mk-def:IsAlbanese, mk-thm:albanese_universal_property_northstar}
  Let \(C/k\) be a smooth proper geometrically integral curve whose Picard functor
  is representable. Then for every \(P\in C(k)\), the scheme \(\Pic^0_{C/k}\)
  satisfies the Albanese universal property for the pointed curve \((C,P)\).
\end{lemma}

\begin{proof}
  \uses{mk-lem:pic0_isAlbanese_algClosed}
  For \(g(C) > 0\) and \(k\) algebraically closed this is
  \ref{mk-lem:pic0_isAlbanese_algClosed}.

  Three extensions are required. Over a general \(k\), apply the algebraically
  closed case over \(k^s\) and descend: the universal morphism and the factoring
  homomorphism are unique, so they are Galois-equivariant and descend to \(k\).
  For \(g(C) = 0\) the tangent space \(H^1(C,\mathcal O_C)\) vanishes, so
  \(\Pic^0_{C/k}\) is a smooth proper group scheme of dimension \(0\), namely
  \(\Spec k\); the universal property then asserts that every pointed morphism
  from \(C\) to an abelian variety is constant, which is the rigidity statement
  that a morphism from a genus-\(0\) curve to an abelian variety is constant.
  The third is the basepoint condition \(P \circ \iota_P = \eta\), a hypothesis
  of \ref{mk-lem:pic0_isAlbanese_algClosed} rather than a restriction on it.
\end{proof}

\begin{lemma}[Albanese property over an algebraically closed field]
  \label{mk-lem:pic0_isAlbanese_algClosed}
  \lean{AlgebraicGeometry.isAlbanese_pic0_of_isAlgClosed}
  \dcref{custom}

  \uses{mk-def:IsAlbanese, mk-thm:albanese_universal_property_northstar}
  Let \(C\) be a smooth proper geometrically integral curve over an
  algebraically closed field \(\bar k\), with representable Picard functor and
  genus \(g(C) > 0\), and let \(P \in C(\bar k)\) satisfy
  \(P \circ \iota_P = \eta\) for the Abel--Jacobi morphism \(\iota_P\) of
  \((C, P)\). Then \(\Pic^0_{C/\bar k}\) satisfies the Albanese universal
  property for the pointed curve \((C, P)\).
\end{lemma}

\begin{proof}
  \uses{mk-thm:albanese_universal_property_northstar}
  The universal property of \ref{mk-def:IsAlbanese} has two clauses. The basepoint
  clause is the hypothesis. The factorisation clause --- every pointed morphism
  \(\varphi : C \to A\) into an abelian variety factors uniquely through
  \(\iota_P\) --- is \ref{mk-thm:albanese_universal_property_northstar} at the same
  object, the Albanese development's Jacobian scheme being \(\Pic^0_{C/\bar k}\).
\end{proof}

\begin{theorem}[Existence of a Jacobian witness]
  \label{mk-thm:nonempty_jacobianWitness}
  \lean{AlgebraicGeometry.nonempty_jacobianWitness}
  \dcref{custom}
  \uses{mk-def:JacobianWitness, mk-def:picardJacobianWitness}

  Every smooth proper geometrically irreducible curve has a Jacobian witness.
\end{theorem}

\begin{proof}
  Take the witness on \(\Pic^0_{C/k}\) supplied by
  Theorem~\ref{mk-def:picardJacobianWitness}.
\end{proof}

\begin{definition}[Chosen Jacobian witness]
  \label{mk-def:jacobianWitness}
  \lean{AlgebraicGeometry.jacobianWitness}
  \dcref{custom}
  \uses{mk-def:JacobianWitness, mk-thm:nonempty_jacobianWitness}

  Fix one Jacobian witness for \(C\).
\end{definition}

\section{The Jacobian and its structure}

\begin{definition}[Jacobian]
  \label{mk-def:Jacobian}
  \lean{AlgebraicGeometry.Jacobian}
  \dcref{custom}
  \uses{mk-def:genus, mk-def:JacobianWitness, mk-def:jacobianWitness}

  The \emph{Jacobian} \(\Jac(C)\) is the underlying \(k\)-scheme of the
  chosen Jacobian witness.  Equivalently, \(\Jac(C)=\Pic^0_{C/k}\) up to the
  unique isomorphism of Albanese objects.
\end{definition}

\begin{theorem}[Group structure]
  \label{mk-thm:Jacobian_grpObj}
  \lean{AlgebraicGeometry.Jacobian.instGrpObj}
  \dcref{custom}
  \uses{mk-def:Jacobian, mk-def:jacobianWitness}

  The Jacobian is a group scheme over \(k\).
\end{theorem}

\begin{proof}
  This is the group structure carried by the chosen Jacobian witness.
\end{proof}

\begin{theorem}[Dimension and smoothness]
  \label{mk-thm:Jacobian_smooth_genus}
  \lean{AlgebraicGeometry.Jacobian.smoothOfRelativeDimension_genus}
  \dcref{custom}
  \uses{mk-def:Jacobian, mk-def:genus, mk-def:jacobianWitness}

  The structural morphism \(\Jac(C)\to\Spec k\) is smooth of relative
  dimension \(g(C)\).
\end{theorem}

\begin{proof}
  This is the dimension assertion carried by the chosen Jacobian witness.
\end{proof}

\begin{theorem}[Properness]
  \label{mk-thm:Jacobian_proper}
  \lean{AlgebraicGeometry.Jacobian.instIsProper}
  \dcref{custom}
  \uses{mk-def:Jacobian, mk-def:jacobianWitness}

  The structural morphism \(\Jac(C)\to\Spec k\) is proper.
\end{theorem}

\begin{proof}
  This is the properness assertion carried by the chosen Jacobian witness.
\end{proof}

\begin{theorem}[Geometric irreducibility]
  \label{mk-thm:Jacobian_geomIrred}
  \lean{AlgebraicGeometry.Jacobian.instGeometricallyIrreducible}
  \dcref{custom}
  \uses{mk-def:Jacobian, mk-def:jacobianWitness}

  The structural morphism \(\Jac(C)\to\Spec k\) is geometrically irreducible.
\end{theorem}

\begin{proof}
  This is the geometric-irreducibility assertion carried by the chosen
  Jacobian witness.
\end{proof}

\begin{remark}
  If \(g(C)=0\), then \(\Jac(C)=\Spec k\).  If \(g(C)=1\) and \(C(k)\) is
  nonempty, the choice of a rational point identifies \(C\) with its
  Jacobian.  For \(g(C)\geq2\), the Jacobian has dimension \(g(C)\), whereas
  \(C\) has dimension one.
\end{remark}

\section{The Albanese property of \(\Pic^0\)}
\label{mk-sec:Jacobian_routeA4_albaneseUP}

\begin{theorem}
[Albanese universal property of \(\Pic^0_{C/k}\)]
  \label{mk-thm:albanese_universal_property_northstar}
  \dcref{custom}
  \uses{mk-lem:abel_jacobi_morphism, mk-lem:symmetric_product_av_map,
    mk-lem:symmetric_product_to_jacobian, mk-lem:descent_through_birational_sigma,
    mk-lem:albanese_eq_iff_symmetricPower_eq, mk-thm:rational_map_to_av_extends,
    mk-thm:rigidity_lemma, mk-lem:av_regular_map_is_hom,
    mk-lem:hom_additivity_over_product}
  Fix \(P\in C(k)\), put \(J=\Pic^0_{C/k}\), and let
  \(f^P(Q)=[\mathcal O_C(Q-P)]\).  For every pointed morphism
  \(\varphi\colon C\to A\) to an abelian variety there is a unique
  group-scheme homomorphism \(\psi\colon J\to A\) such that
  \(\varphi=f^P\circ\psi\).
\end{theorem}

\begin{proof}
  For \(g=g(C)>0\), the morphism
  \[
    (Q_1,\ldots,Q_g)\longmapsto
      \varphi(Q_1)+\cdots+\varphi(Q_g)
  \]
  is symmetric and therefore descends to \(C^{(g)}\).  The Abel map
  \(C^{(g)}\to J\) is birational onto \(J\), so this descent defines a
  rational map \(J\dashrightarrow A\).  A rational map from a nonsingular
  variety to an abelian variety extends uniquely to a regular morphism;
  denote the extension by \(\psi\).  By construction
  \(\varphi=f^P\circ\psi\), and \(\psi(0)=0\), so rigidity makes \(\psi\) a
  homomorphism.  If \(\psi'\) is another such homomorphism, then \(\psi\) and
  \(\psi'\) agree on the sum of \(g\) copies of \(f^P(C)\), which is all of
  \(J\); hence \(\psi=\psi'\).  When \(g=0\), \(J=\Spec k\) and the result
  follows from rigidity.
\end{proof}
